\section{Future Work Beyond the Next Paper}
\label{sec:futurework}

\subsection{LLM Agents as Substitutes for Human Participants}

An important direction for future work is to evaluate whether LLM-based agents can reliably simulate human code comprehension behavior. Building on our findings that time-based and input--output reasoning proxies align closely with expert judgments (Section~\ref{ase}), we propose to replicate our human study protocol using LLM agents as participants. Specifically, we will measure whether the relative performance of LLMs on comprehension tasks (e.g., response time to determine outputs given inputs) exhibits similar patterns to expert-consensus rankings. Beyond replicating our own study, we will also attempt to replicate selected prior comprehensibility studies that used human participants, assessing how well LLMs reproduce the patterns observed in those studies. This broadens the contribution from validating a single protocol to assessing the general viability of LLM-as-participant as a methodology in comprehensibility research. This will allow us to assess the extent to which LLMs can serve as scalable substitutes for human participants in controlled comprehension studies, potentially reducing the cost and effort associated with human-subject experiments.

\subsection{Reassessment of Prior Comprehensibility Studies}

Another promising direction is to conduct a comprehensive reassessment of prior code comprehensibility studies in light of our findings on proxy reliability. A large portion of the literature relies on proxies that our results suggest are weak or even negatively correlated with actual comprehension difficulty (e.g.,~\cite{Wyrich:ACS23,33-29,33-4}), while some studies combine such unreliable proxies with more reliable ones (e.g.,~\cite{33-15,penningtonModel,S75,S4}). At the same time, other studies employ proxies that appear to be more strongly aligned with expert-consensus judgments (e.g.,~\cite{S39,S48,S50,S57,S59,S60}). By systematically revisiting and reinterpreting these studies through the lens of validated, reliable proxies, we aim to reassess their conclusions and determine which findings remain robust and which may require revision. This effort will provide a clearer understanding of the reliability of existing results in the literature and contribute to establishing stronger methodological foundations for future research.

\subsection{Accidental Complexity and Code Comprehensibility}

Another direction is to investigate the relationship between accidental complexity~\cite{brooks1987no} and code comprehensibility, and in particular to examine the relative impact of different kinds of accidental complexity on comprehension difficulty. Building on our ground-truth dataset, we will systematically introduce controlled amounts of accidental complexity into code snippets through semantics-preserving transformations (e.g., adding redundant control flow, unnecessary abstractions, or convoluted naming). Rather than rerunning the expert study, we will rely on input--output response time as the outcome measure, which our ASE findings validate as a reliable proxy for comprehension difficulty. We will analyze how different transformation types affect comprehension time, allowing us to assess which forms of accidental complexity most strongly hinder human understanding. This provides deeper insight into the limits of current proxies and guides the design of more robust measures of code comprehensibility.

\subsection{Impact of Java Constructs on Comprehensibility}

Another direction is to investigate how different code constructs in Java affect comprehensibility using the reliable proxies identified in our ASE study. Prior work (e.g.,~\cite{lucas2019does}) has examined the impact of constructs such as control flow patterns or API usage, often relying on proxies whose reliability is unclear or limited, although some studies employ more reliable measures (e.g.,~\cite{mehlhorn2022imperative}). By leveraging the proxies validated in our work (Section~\ref{ase}), we will systematically reevaluate these constructs and, where prior studies relied on weak proxies, replicate and reassess their findings using more reliable measures. This approach allows us to identify which patterns consistently hinder or facilitate comprehension, and to confirm, refine, or revise conclusions drawn in earlier work. Ultimately, this line of research can provide deeper insights into how Java language features and coding idioms influence human understanding of code.

\subsection{Alternative Automatic Proxies for Verifiability}

Another direction is to identify and evaluate alternative automatic proxies for code verifiability. In our current work, verifiability is operationalized using raw warning counts produced by selected Checker Framework checkers, treating all warnings as false positives. We propose to explore additional candidate proxies that approximate verification effort through different signals, and to validate them through controlled human studies. Specifically, we will assess which proxies most reliably reflect the underlying construct of verifiability, defined as the effort required to establish program correctness using verification tools. Once robust proxies for verifiability are identified, we will further examine their relationship with code comprehensibility to determine whether stronger or more consistent effects can be established.

\subsection{Reducing False Positives via Comprehensibility Insights}

If a strong relationship between code verifiability and comprehensibility is found (Section~\ref{sec:nextpaper}), we propose to investigate techniques for reducing false positive warnings produced by verification tools based on comprehensibility insights. Many verification tools produce a large number of false positives due to the conservatism required for soundness. By leveraging insights about code regions that are harder to understand, we can explore whether such regions are more prone to false positives and design approaches to mitigate them. This may include prioritizing warnings, filtering spurious alerts, or guiding code transformations that both improve comprehensibility and reduce verification noise.

\subsection{Tool Support for Comprehensibility-Driven Refactoring}

Finally, we propose to explore the development of practical tools if a sufficiently reliable automatic proxy for code comprehensibility can be identified. In that case, we will design and evaluate a tool that identifies hard-to-understand regions of code and refactors them into semantically equivalent but more comprehensible versions. Such a tool could use the proxy to detect patterns associated with high comprehension difficulty and apply semantics-preserving transformations to improve readability and understanding. We will evaluate this approach by measuring improvements in the proxy and, where feasible, through targeted human studies to validate that the refactored code is indeed easier to understand. This direction connects the measurement of comprehensibility to actionable improvements in software quality.